% Zone „Das kennst du schon“ – aus bank/prozentrechnung/zone.jsonl (zusammenbau v0.1)
\einheitenkopf[zone][Kennst du schon]{Das kennst du schon}
\setcounter{aufgabe}{0}
%% TODO Grundvorstellungs-Aufgabe als erste Hauptnummer der Zone (2.8) fehlt in der Bank

%% TODO Titel: Ich-kann-Satz fehlt in der Bank (Platzhalter: Kettenname) – Nr. 1
%% TODO Anweisung: ein Satz über der ersten Teilaufgabe fehlt in der Bank – Nr. 1
\begin{aufgabe}{Bruch als Anteil lesen und kürzen (drei von zwölf, gekürzt ein Viertel), auf den Nenner hundert erweitern}
%% TODO Darstellung für schwach fehlt in der Bank (prozentrechnung-zone-f1-v1; Abschnitt „Für schwache Schüler“)
\swz{$6$ von $10$ Bonbons sind rot – welcher Anteil? Kürze.}{}
%% TODO Darstellung für schwach fehlt in der Bank (prozentrechnung-zone-f1-v3; Abschnitt „Für schwache Schüler“)
\swz{$6$ von $15$ Plätzen sind besetzt – welcher Anteil? Kürze und erweitere auf Hundertstel.}{}
\end{aufgabe}

%% TODO Titel: Ich-kann-Satz fehlt in der Bank (Platzhalter: Kettenname) – Nr. 2
%% TODO Anweisung: ein Satz über der ersten Teilaufgabe fehlt in der Bank – Nr. 2
\begin{aufgabe}{Bruch ↔ Dezimalzahl (ein Viertel, drei Fünftel)}
%% TODO Darstellung für schwach fehlt in der Bank (prozentrechnung-zone-f2-v1; Abschnitt „Für schwache Schüler“)
\swz{$\frac{1}{2}$ als Dezimalzahl?}{}
%% TODO Darstellung für schwach fehlt in der Bank (prozentrechnung-zone-f2-v3; Abschnitt „Für schwache Schüler“)
\swz{$\frac{3}{25}$ als Dezimalzahl?}{}
\end{aufgabe}

%% TODO Titel: Ich-kann-Satz fehlt in der Bank (Platzhalter: Kettenname) – Nr. 3
%% TODO Anweisung: ein Satz über der ersten Teilaufgabe fehlt in der Bank – Nr. 3
\begin{aufgabe}{Durch hundert teilen, mit Dezimalzahlen multiplizieren (Kommaverschiebung)}
%% TODO Darstellung für schwach fehlt in der Bank (prozentrechnung-zone-f3-v1; Abschnitt „Für schwache Schüler“)
\swz{$600 : 100$}{}
%% TODO Darstellung für schwach fehlt in der Bank (prozentrechnung-zone-f3-v3; Abschnitt „Für schwache Schüler“)
\swz{$0{,}15 \cdot 60$}{}
\end{aufgabe}

%% TODO Titel: Ich-kann-Satz fehlt in der Bank (Platzhalter: Kettenname) – Nr. 4
%% TODO Anweisung: ein Satz über der ersten Teilaufgabe fehlt in der Bank – Nr. 4
\begin{aufgabe}{Hoch- und Runterrechnen in einer Tabelle (Dreisatz, proportionale Zuordnung)}
\swz{$2$ Kinokarten kosten $18\,€$ – wie viel kosten $5$ Karten?}{\begin{dreisatz}{Karten}{Preis}\dsz{2}{18\,€}\dsp{\phantom{:0}}{\phantom{:0}}\dsz{1}{\dsleer}\dsp{\phantom{:0}}{\phantom{:0}}\dsz{5}{\dsleer}\end{dreisatz}}
\swz{$5$ Brötchen kosten $2\,€$ – wie viel kosten $3$ Brötchen?}{\begin{dreisatz}{Brötchen}{Preis}\dsz{5}{2\,€}\dsp{\phantom{:0}}{\phantom{:0}}\dsz{1}{\dsleer}\dsp{\phantom{:0}}{\phantom{:0}}\dsz{3}{\dsleer}\end{dreisatz}}
\end{aufgabe}

%% TODO Titel: Ich-kann-Satz fehlt in der Bank (Platzhalter: Kettenname) – Nr. 5
%% TODO Anweisung: ein Satz über der ersten Teilaufgabe fehlt in der Bank – Nr. 5
\begin{aufgabe}{Runden auf eine Dezimalstelle}
%% TODO Darstellung für schwach fehlt in der Bank (prozentrechnung-zone-f5-v1; Abschnitt „Für schwache Schüler“)
\swz{$4{,}73$ – auf eine Dezimalstelle?}{}
%% TODO Darstellung für schwach fehlt in der Bank (prozentrechnung-zone-f5-v3; Abschnitt „Für schwache Schüler“)
\swz{$18{,}456$ kg – auf eine Dezimalstelle?}{}
\end{aufgabe}

%% TODO Titel: Ich-kann-Satz fehlt in der Bank (Platzhalter: Kettenname) – Nr. 6
%% TODO Anweisung: ein Satz über der ersten Teilaufgabe fehlt in der Bank – Nr. 6
\begin{aufgabe}{Bruchteil einer Größe (zwei Drittel von 60 €)}
%% TODO Darstellung für schwach fehlt in der Bank (prozentrechnung-zone-f6-v1; Abschnitt „Für schwache Schüler“)
\swz{$\frac{1}{4}$ von $20\,€$}{}
%% TODO Darstellung für schwach fehlt in der Bank (prozentrechnung-zone-f6-v3; Abschnitt „Für schwache Schüler“)
\swz{$\frac{3}{5}$ von $45$ m}{}
\end{aufgabe}

%% TODO Titel: Ich-kann-Satz fehlt in der Bank (Platzhalter: Kettenname) – Nr. 7
%% TODO Anweisung: ein Satz über der ersten Teilaufgabe fehlt in der Bank – Nr. 7
\begin{aufgabe}{Bruchteil einer Größe (zwei Drittel von 60 €) – Fehler finden}
\swz[3]{Von $40\,€$ ist ein Fünftel ausgegeben – wie viel ist übrig? Ben rechnet: $40 : 5 = 8$, übrig sind $8\,€$. Wo ist der Fehler?}{}
\end{aufgabe}

%% TODO Titel: Ich-kann-Satz fehlt in der Bank (Platzhalter: Kettenname) – Nr. 8
%% TODO Anweisung: ein Satz über der ersten Teilaufgabe fehlt in der Bank – Nr. 8
\begin{aufgabe}{Bruchteil einer Größe (zwei Drittel von 60 €) – selbst rechnen}
%% TODO Darstellung für schwach fehlt in der Bank (prozentrechnung-zone-f6-v7; Abschnitt „Für schwache Schüler“)
\swz{Von $30$ Litern im Tank ist ein Drittel verbraucht – wie viel ist übrig?}{}
\end{aufgabe}
